% Source: Weinberg GR, physical PDF pp. 69--72
%         (printed pp. 43--46).
% Coverage: Section 2.8, Energy-Momentum Tensor; equations
%           (2.8.1)--(2.8.14), including (2.8.5a).
% Figures/tables/footnotes: none.
% Status: source-reviewed and compile-clean.
% Uncertainties: none.

\subsection{Energy-Momentum Tensor}\label{sec:2.8}

In Section~\ref{sec:2.6} we introduced the density \(\rho\) and current
\(\mathbf J\) of electric charge. We now give a similar definition for the
density and current of the energy-momentum four-vector \(p^\mu\). First
consider a system of particles labeled \(n\), with energy-momentum four-vectors
\(p_n^\mu(t)\). The density of \(p^\mu\) is defined by
\begin{equation}
  T^{\mu0}(\mathbf x,t)
  \equiv\sum_n p_n^\mu(t)
  \delta^3\bigl(\mathbf x-\mathbf x_n(t)\bigr),
  \tag{2.8.1}\label{eq:2.8.1}
\end{equation}
and its current is defined by
\begin{equation}
  T^{\mu i}(\mathbf x,t)
  \equiv\sum_n p_n^\mu(t)\frac{dx_n^i(t)}{dt}
  \delta^3\bigl(\mathbf x-\mathbf x_n(t)\bigr).
  \tag{2.8.2}\label{eq:2.8.2}
\end{equation}
These two definitions can be united into a single formula,
\begin{equation}
  T^{\mu\nu}(x)
  =\sum_n p_n^\mu(t)\frac{dx_n^\nu(t)}{dt}
  \delta^3\bigl(\mathbf x-\mathbf x_n(t)\bigr),
  \tag{2.8.3}\label{eq:2.8.3}
\end{equation}
where \(x_n^0(t)=t\). We note from Eq.~\eqref{eq:2.4.10} that
\[
  p_n^\nu=E_n\frac{dx_n^\nu}{dt},
\]
so Eq.~\eqref{eq:2.8.3} can also be written as
\begin{equation}
  T^{\mu\nu}(x)
  =\sum_n\frac{p_n^\mu p_n^\nu}{E_n}
  \delta^3\bigl(\mathbf x-\mathbf x_n(t)\bigr),
  \tag{2.8.4}\label{eq:2.8.4}
\end{equation}
and we see that \(T^{\mu\nu}\) is symmetric:
\begin{equation}
  T^{\mu\nu}(x)=T^{\nu\mu}(x).
  \tag{2.8.5}\label{eq:2.8.5}
\end{equation}
We can also write Eq.~\eqref{eq:2.8.3}, in analogy with
Eq.~\eqref{eq:2.6.5}, as
\begin{equation}
  T^{\mu\nu}(x)
  =\sum_n\int d\tau\,p_n^\mu
    \frac{dx_n^\nu}{d\tau}\,
    \delta^4\bigl(x-x_n(\tau)\bigr),
  \tag{2.8.5a}\label{eq:2.8.5a}
\end{equation}
and we see that \(T^{\mu\nu}\) is a tensor; that is,
\[
  T'^{\mu\nu}
  =\Lambda^\mu{}_\rho\Lambda^\nu{}_\sigma T^{\rho\sigma}
\]
under a Lorentz transformation \eqref{eq:2.1.1}.

The conservation law for \(T^{\mu\nu}\) will take a little more thought.
Returning to Eqs.~\eqref{eq:2.8.1} and \eqref{eq:2.8.2}, we see that
\[
\begin{aligned}
  \partial_i T^{\mu i}(\mathbf x,t)
  &=-\sum_n p_n^\mu(t)\frac{dx_n^i(t)}{dt}
    % NOTATION EXCEPTION: the differentiated particle coordinate must remain
    % explicit to distinguish it from the field coordinate.
    \frac{\partial}{\partial x_n^i}
    \delta^3\bigl(\mathbf x-\mathbf x_n(t)\bigr)\\
  &=-\sum_n p_n^\mu(t)\partial_t
    \delta^3\bigl(\mathbf x-\mathbf x_n(t)\bigr)\\
  &=-\partial_t T^{\mu0}(\mathbf x,t)
    +\sum_n\frac{dp_n^\mu(t)}{dt}
     \delta^3\bigl(\mathbf x-\mathbf x_n(t)\bigr),
\end{aligned}
\]
and so
\begin{equation}
  \partial_\nu T^{\mu\nu}=G^\mu ,
  \tag{2.8.6}\label{eq:2.8.6}
\end{equation}
where \(G^\mu\) is the \emph{density of force},
\[
\begin{aligned}
  G^\mu(\mathbf x,t)
  &\equiv\sum_n\delta^3\bigl(\mathbf x-\mathbf x_n(t)\bigr)
      \frac{dp_n^\mu(t)}{dt}\\
  &=\sum_n\delta^3\bigl(\mathbf x-\mathbf x_n(t)\bigr)
      \frac{d\tau}{dt}f_n^\mu(t).
\end{aligned}
\]
If the particles are free, then \(p_n^\mu\) is constant and \(T^{\mu\nu}\)
is conserved; that is,
\begin{equation}
  \partial_\nu T^{\mu\nu}(x)=0 .
  \tag{2.8.7}\label{eq:2.8.7}
\end{equation}
The same is also true if the particles interact only during collisions that are
strictly localized in space. In this case Eq.~\eqref{eq:2.8.6} gives
\[
  \partial_\nu T^{\mu\nu}(x)
  =\sum_c\delta^3\bigl(\mathbf x-\mathbf x_c(t)\bigr)
    \frac{d}{dt}\sum_{n\in c}p_n^\mu(t),
\]
where \(\mathbf x_c(t)\) is the location of the \(c\)th collision going on at
time \(t\), and \(n\in c\) means we sum only over the particles participating
in the \(c\)th collision. But each collision conserves momentum, so
\(\sum_{n\in c}p_n^\mu(t)\) must be time-independent, yielding the
conservation equation \eqref{eq:2.8.7}.

The energy-momentum tensor \eqref{eq:2.8.3} will not be conserved if the
particles are subject to forces that act at a distance. For instance, consider
a gas of charged particles, with charges \(e_n\). Then
Eqs.~\eqref{eq:2.8.6}, \eqref{eq:2.4.1}, and \eqref{eq:2.7.9} give
\[
  \partial_\nu T^{\mu\nu}(x)
  =\sum_n e_nF^\mu{}_\rho(x)\frac{dx_n^\rho}{dt}
    \delta^3\bigl(\mathbf x-\mathbf x_n(t)\bigr),
\]
and, using Eq.~\eqref{eq:2.6.4}, this gives
\begin{equation}
  \partial_\nu T^{\mu\nu}(x)
  =F^\mu{}_\rho(x)J^\rho(x).
  \tag{2.8.8}\label{eq:2.8.8}
\end{equation}
Although this is not conserved, we can construct a conserved tensor by adding
a purely electromagnetic term
\begin{equation}
  T_{\mathrm{em}}^{\mu\nu}
  \equiv F^\mu{}_\rho F^{\nu\rho}
  -\frac14\eta^{\mu\nu}F_{\rho\sigma}F^{\rho\sigma}.
  \tag{2.8.9}\label{eq:2.8.9}
\end{equation}
That is, the electromagnetic energy and momentum densities are given by
\begin{equation}
  T_{\mathrm{em}}^{00}
  =\frac12(\mathbf E^2+\mathbf B^2),
  \qquad
  T_{\mathrm{em}}^{i0}=(\mathbf E\times\mathbf B)_i .
  \tag{2.8.10}\label{eq:2.8.10}
\end{equation}
We note that
\[
  \partial_\nu T_{\mathrm{em}}^{\mu\nu}
  =F^\mu{}_\rho\partial_\nu F^{\nu\rho}
   +F^{\nu\rho}\partial_\nu F^\mu{}_\rho
   -\frac12F_{\rho\sigma}\partial^\mu F^{\rho\sigma}.
\]
[Here \(\partial^\mu=\eta^{\mu\nu}\partial_\nu\).] With a little reshuffling
of indices, this becomes
\[
  \partial_\nu T_{\mathrm{em}}^{\mu\nu}
  =F^\mu{}_\rho\partial_\nu F^{\nu\rho}
   -\frac12F^{\rho\sigma}
   \left(
     \partial^\mu F_{\rho\sigma}
     +\partial_\rho F_\sigma{}^\mu
     +\partial_\sigma F^\mu{}_\rho
   \right).
\]
Using the Maxwell equations \eqref{eq:2.7.6} and \eqref{eq:2.7.10}, we find
\begin{equation}
  \partial_\nu T_{\mathrm{em}}^{\mu\nu}
  =-F^\mu{}_\rho J^\rho .
  \tag{2.8.11}\label{eq:2.8.11}
\end{equation}
Comparing Eq.~\eqref{eq:2.8.8} with Eq.~\eqref{eq:2.8.11}, we are led to
redefine the energy-momentum tensor as
\begin{equation}
  T^{\mu\nu}
  =\sum_n p_n^\mu\frac{dx_n^\nu}{dt}
    \delta^3\bigl(\mathbf x-\mathbf x_n(t)\bigr)
   +T_{\mathrm{em}}^{\mu\nu}.
  \tag{2.8.12}\label{eq:2.8.12}
\end{equation}
This is again a symmetric tensor, and is now conserved:
\begin{equation}
  \partial_\nu T^{\mu\nu}=0 .
  \tag{2.8.13}\label{eq:2.8.13}
\end{equation}
We can continue to add more and more terms to \(T^{\mu\nu}\) to account for
other fields and keep \(T^{\mu\nu}\) conserved. A systematic method for
constructing these terms is presented in Chapter~12.

Just as the integral of the charge density \(J^0\) is the total charge, the
integral of the density \(T^{\mu0}\) of \(p^\mu\) is the total \(p^\mu\):
\begin{equation}
  p^\mu_{\mathrm{total}}
  =\int d^3x\,T^{\mu0}(\mathbf x,t).
  \tag{2.8.14}\label{eq:2.8.14}
\end{equation}
That this is a constant four-vector can be shown in the same way that we showed
in Section~\ref{sec:2.6} that the total charge \eqref{eq:2.6.7} is a constant
scalar.
